\documentclass[../../../../main.tex]{subfiles}
\begin{document}

$\theta_1$，$\theta_2$，$\theta_3$，$\theta_4$，$\theta_5$ を正の数とする．
図のように円に内接する5角形 $A_1A_2A_3A_4A_5$ で，
$1 \leqq i \leqq 5$ に対し角 $A_i$ の大きさが $\theta_i$ となるものが存在するためには，
\begin{align*}
  \theta_1+\theta_2+\theta_3+\theta_4+\theta_5=3\pi, \hspace{1zw}
  \theta_1+\theta_3>\pi, \hspace{1zw}
  \theta_2+\theta_4>\pi, \hspace{1zw}
  \theta_3+\theta_5>\pi, \hspace{1zw}
  \theta_1+\theta_4>\pi, \hspace{1zw}
  \theta_2+\theta_5>\pi
\end{align*}
が同時に成り立つことが必要かつ十分であることを証明せよ．

\begin{figure}[htbp]
  \centering
  \begin{tikzpicture}[scale=1.5]
    \draw[thick] (0,0) circle (1.5);

    \coordinate (A1) at (-25:1.5);
    \coordinate (A2) at (40:1.5);
    \coordinate (A3) at (125:1.5);
    \coordinate (A4) at (195:1.5);
    \coordinate (A5) at (260:1.5);

    \draw[thick] (A1) -- (A2) -- (A3) -- (A4) -- (A5) -- cycle;

    \fill (A1) circle (1.2pt) node[below right] {$A_1$};
    \fill (A2) circle (1.2pt) node[above right] {$A_2$};
    \fill (A3) circle (1.2pt) node[above left] {$A_3$};
    \fill (A4) circle (1.2pt) node[left] {$A_4$};
    \fill (A5) circle (1.2pt) node[below] {$A_5$};

    \draw (A1) +(105:0.35) arc (105:202:0.35);
    \node at ($(A1)+(150:0.5)$) {$\theta_1$};

    \draw (A2) +(195:0.35) arc (195:262:0.35);
    \node at ($(A2)+(230:0.5)$) {$\theta_2$};

    \draw (A3) +(-75:0.35) arc (-75:15:0.35);
    \node at ($(A3)+(-30:0.5)$) {$\theta_3$};

    \draw (A4) +(15:0.35) arc (15:60:0.35);
    \node at ($(A4)+(35:0.5)$) {$\theta_4$};

    \draw (A5) +(80:0.35) arc (80:150:0.35);
    \node at ($(A5)+(115:0.5)$) {$\theta_5$};
  \end{tikzpicture}
\end{figure}

\end{document}
